\chapter{Gaussian processes and the Brownian bridge}
\label{ch:gaussian}

\begin{question}
Brownian motion at many times at once: what is the joint law, and which processes built from Brownian motion are still ``normal everywhere''?
\end{question}

\begin{whatfor}
\begin{itemize}
  \item Defines the multivariate normal (MVN) distribution, its covariance matrix and joint exponential moment.
  \item Defines Gaussian processes and shows that Brownian motion is the Gaussian process with mean $0$ and covariance $\min(s,t)$.
  \item Studies two Gaussian processes built from it: the Brownian bridge and integrated Brownian motion.
\end{itemize}
\skipline{fs:gaussian}
\end{whatfor}

\begin{prereq}
From Part~0: covariance and its rules, independence, sums of independent normals, the normal exponential moment, CE1--CE4, matrices. From Chapters~6, 8, 10: Wiener process, independent increments, conditional law of $W(s)$ given $W(t)$.
\end{prereq}

%=====================================================================
\section{The multivariate normal distribution}

\begin{definition}[Multivariate normal]
\label{def:mvn}
Let $Z_1,\dots,Z_n$ be independent $N(0,1)$. Random variables $X_1,\dots,X_m$ have a \dfn{multivariate normal} (MVN) distribution if each is a fixed linear combination of the $Z_j$ plus a constant:
\[
  X_i=a_{i1}Z_1+\dots+a_{in}Z_n+\mu_i,\qquad i=1,\dots,m;\qquad\text{in matrix form } X=AZ+\mu .
\]
\end{definition}

The numbers that matter are the means and covariances:
\[
  \E_P(X_i)=\mu_i,\qquad \Cov(X_i,X_j)=\sum_{k=1}^na_{ik}a_{jk},\qquad\text{i.e.\ the covariance matrix is }\Sigma=AA^{\mathsf T}.
\]

\begin{result}[Linear combinations and the joint exponential moment]
\label{res:mvn}
If $X=(X_1,\dots,X_m)$ is MVN with means $\mu_i$ and covariances $\Cov(X_i,X_j)$, then for any numbers $t_1,\dots,t_m$:
\begin{enumerate}[label=(\alph*)]
  \item $\sum_it_iX_i$ is normal, with mean $\sum_it_i\mu_i$ and variance $\sum_i\sum_jt_it_j\Cov(X_i,X_j)$;
  \item the \dfn{joint exponential moment} is
  \[
    \E_P\,e^{t_1X_1+\dots+t_mX_m}=\exp\Big\{\sum_it_i\mu_i+\frac12\sum_i\sum_jt_it_j\Cov(X_i,X_j)\Big\}.
  \]
\end{enumerate}
So an MVN distribution is completely determined by its means and covariances.
\end{result}
\begin{proof}
Everything reduces to one normal variable.
\begin{steps}
  \item $\sum_it_iX_i=\sum_j\big(\sum_it_ia_{ij}\big)Z_j+\sum_it_i\mu_i$ is a constant plus a combination of independent normals, hence normal (\cref{res:normal-sum}). Its mean and variance follow from the rules of Chapter~3.
  \item For (b), apply $\E_Pe^{Y}=e^{\E_PY+\frac12\Var Y}$ to $Y=\sum_it_iX_i$.
  \item The joint exponential moment depends only on the means and covariances, and (as in Chapter~4) it determines the joint distribution. \qedhere
\end{steps}
\end{proof}

\begin{result}[Uncorrelated means independent, for MVN]
\label{res:uncorr}
If $(X,Y)$ is MVN and $\Cov(X,Y)=0$, then $X$ and $Y$ are independent.
\end{result}
\begin{proof}
With $\Cov=0$ the double sum in \cref{res:mvn}(b) splits: $\E_Pe^{t_1X+t_2Y}=e^{t_1\mu_X+\frac12t_1^2\Var X}\cdot e^{t_2\mu_Y+\frac12t_2^2\Var Y}=\E_Pe^{t_1X}\,\E_Pe^{t_2Y}$. This is the joint exponential moment of an independent pair with the same marginals, so the joint law is that of an independent pair.
\end{proof}

This is special to MVN. In general, zero covariance does not give independence.

\begin{result}[Conditioning within an MVN pair]
\label{res:mvn-cond}
If $(X,Y)$ is MVN, then $X$ given $Y=y$ is normal with
\[
  \E_P(X\given Y=y)=\mu_X+\frac{\Cov(X,Y)}{\Var Y}(y-\mu_Y),\qquad \Var(X\given Y=y)=\Var X-\frac{\Cov(X,Y)^2}{\Var Y}.
\]
\end{result}
\begin{proof}
Let $c=\Cov(X,Y)/\Var Y$ and $R=X-cY$. Then $\Cov(R,Y)=\Cov(X,Y)-c\Var Y=0$, and $(R,Y)$ is MVN, so $R$ is independent of $Y$. Given $Y=y$, $X=cy+R$ with $R\sim N(\mu_X-c\mu_Y,\ \Var X-c^2\Var Y)$.
\end{proof}

For $(W(s),W(t))$ with $s<t$ this gives back Chapter~9: $\Cov=s$, $\Var W(t)=t$, so the mean is $\frac stW(t)$ and the variance $s-\frac{s^2}{t}=\frac{s(t-s)}{t}$.

%=====================================================================
\section{Gaussian processes}

\begin{definition}[Gaussian process]
A process $(X(t))_{t\ge0}$ is a \dfn{Gaussian process} if for all times $t_1,\dots,t_n$, the vector $(X(t_1),\dots,X(t_n))$ is MVN. Its \dfn{mean function} is $m(t)=\E_PX(t)$ and its \dfn{covariance function} is $K(s,t)=\Cov\big(X(s),X(t)\big)$.
\end{definition}

By \cref{res:mvn}, a Gaussian process is described completely by $m$ and $K$.

\begin{theorem}[Brownian motion is Gaussian]
\label{thm:bm-gauss}
Standard Brownian motion is a Gaussian process with $m(t)=0$ and $K(s,t)=\min(s,t)$. Conversely, a Gaussian process with continuous paths, $X(0)=0$, mean $0$ and covariance $\min(s,t)$ is a standard Brownian motion.
\end{theorem}
\begin{proof}
Write values as sums of independent increments.
\begin{steps}
  \item For $t_1<\dots<t_n$: $W(t_k)=W(t_1)+\big(W(t_2)-W(t_1)\big)+\dots+\big(W(t_k)-W(t_{k-1})\big)$, a linear combination of independent normals. So the vector is MVN.
  \item For $s<t$: $\Cov(W(s),W(t))=\E_P\big[W(s)\big(W(s)+W(t)-W(s)\big)\big]=\E_PW(s)^2+\E_PW(s)\,\E_P\big(W(t)-W(s)\big)=s+0=\min(s,t)$.
  \item Converse: increments over disjoint intervals $(r,s]$, $(s,t]$ have covariance $\min(s,t)-\min(s,s)-\min(r,t)+\min(r,s)=s-s-r+r=0$, so by \cref{res:uncorr} they are independent; and $\Var\big(X(t)-X(s)\big)=t-2s+s=t-s$. These are the rules of a Wiener process. \qedhere
\end{steps}
\end{proof}

\begin{pic}
\begin{tikzpicture}
\begin{axis}[width=7.2cm,height=6.2cm,view={35}{35},xlabel={$s$},ylabel={$t$},zlabel={$\min(s,t)$},colormap={blu}{color=(defc!15) color=(defc)},samples=21,domain=0:1,y domain=0:1,z buffer=sort]
  \addplot3[surf,shader=faceted interp,draw=defc!40] {min(x,y)};
\end{axis}
\end{tikzpicture}
\picnote{The covariance function $K(s,t)=\min(s,t)$ of Brownian motion: a ridge along the diagonal $s=t$, where it equals the variance $t$.}
\end{pic}

\begin{example}[Two values of Brownian motion]
\label{ex:4w1-6w3}
$4W_1+6W_3$ is normal with mean $0$ and variance $16\cdot1+36\cdot3+2\cdot4\cdot6\cdot\min(1,3)=16+108+48=172$.
\end{example}

\begin{example}[Exponentials of Brownian motion]
\label{ex:cov-exp}
$W_1+W_2\sim N(0,\,1+2+2\cdot1)=N(0,5)$, so $\E_Pe^{W_1+W_2}=e^{5/2}$, and
\[
  \Cov\big(e^{W_1},e^{W_2}\big)=\E_Pe^{W_1+W_2}-\E_Pe^{W_1}\E_Pe^{W_2}=e^{5/2}-e^{1/2}e^{1}=e^{3/2}(e-1)\approx7.70 .
\]
\end{example}

\begin{example}[A time-changed Brownian motion]
\label{ex:time-change}
Is $X(t)=e^{-t}W(e^{2t})$ a Gaussian process? Each $X(t_k)$ is a fixed number times a value of $W$ at a fixed time, so any vector $(X(t_1),\dots,X(t_n))$ is a linear transformation of an MVN vector: yes. Its covariance is
\[
  K(s,t)=e^{-s}e^{-t}\min\big(e^{2s},e^{2t}\big)=e^{-|t-s|} .
\]
It has $X(0)=W(1)\ne0$ and depends only on $|t-s|$: it is a stationary process that returns in Chapter~16 as a model of interest rates, not a Brownian motion.
\end{example}

%=====================================================================
\section{The Brownian bridge}

\begin{definition}[Brownian bridge]
The \dfn{Brownian bridge} is standard Brownian motion on $[0,1]$ conditioned on $W(1)=0$: the process $\{W(t),0\le t\le1\given W(1)=0\}$, tied down at both ends.
\end{definition}

\begin{theorem}[Mean and covariance of the bridge]
\label{thm:bridge}
The Brownian bridge is a Gaussian process with mean $0$ and covariance $K(s,t)=s(1-t)$ for $s\le t$. The process $Z(t)=W(t)-t\,W(1)$, $0\le t\le1$, has the same mean and covariance, so it is a Brownian bridge.
\end{theorem}
\begin{proof}
Use the conditional law of Chapter~9 and the tower rule.
\begin{steps}
  \item Mean: $\E_P\big(W(s)\given W(1)=0\big)=s\cdot0=0$.
  \item Covariance, $s<t<1$: by the tower rule, conditioning first on $W(t)$,
  \[
    \begin{aligned}
    \E_P\big(W(s)W(t)\given W(1)=0\big)&=\E_P\Big[W(t)\,\E_P\big(W(s)\given W(t)\big)\Bigm| W(1)=0\Big]\\
    &=\frac st\,\E_P\big(W(t)^2\given W(1)=0\big)=\frac st\cdot t(1-t)=s(1-t),
    \end{aligned}
  \]
  using $\E_P(W(s)\given W(t),W(1))=\frac stW(t)$ (the value at $1$ adds nothing once $W(t)$ is known) and $\Var(W(t)\given W(1)=0)=t(1-t)$.
  \item For $Z$: it is Gaussian (linear in $W$), $\E_PZ(t)=0$, and for $s<t$
  \[
    \begin{aligned}
    \Cov(Z(s),Z(t))&=\Cov(W(s),W(t))-t\Cov(W(s),W(1))\\
    &\quad-s\Cov(W(1),W(t))+st\Var W(1)\\
    &=s-ts-st+st=s(1-t). \qedhere
    \end{aligned}
  \]
\end{steps}
\end{proof}

\begin{pic}
\begin{tikzpicture}
\begin{axis}[width=10cm,height=4.8cm,xmin=0,xmax=1,xlabel={$t$},ylabel={$W(t)-tW(1)$},samples=100]
  \addplot[draw=none,fill=thmc!12,domain=0:1] {sqrt(x*(1-x))} \closedcycle;
  \addplot[draw=none,fill=thmc!12,domain=0:1] {-sqrt(x*(1-x))} \closedcycle;
  \addplot[defc,thick] table[x=t,y expr=\thisrow{w1}-\thisrow{t}*1.0089] {figs/bm3.dat};
  \addplot[trapc,thick] table[x=t,y expr=\thisrow{w2}+\thisrow{t}*0.8089] {figs/bm3.dat};
  \addplot[prec,thick] table[x=t,y expr=\thisrow{w3}-\thisrow{t}*0.0112] {figs/bm3.dat};
  \addplot[only marks,mark=*] coordinates {(0,0)(1,0)};
\end{axis}
\end{tikzpicture}
\picnote{Three Brownian bridges, made from the Brownian paths of Chapter~6 by subtracting $tW(1)$. All are pinned to $0$ at both ends; the shaded band is $\pm\sqrt{t(1-t)}$, one standard deviation.}
\end{pic}

%=====================================================================
\section{Integrated Brownian motion}

A commodity price whose \emph{rate of change} behaves like Brownian motion is modelled by integrating it.

\begin{definition}[Integrated Brownian motion]
For standard Brownian motion $W$, the process $Z(t)=\int_0^tW(s)\,ds$ is \dfn{integrated Brownian motion}.
\end{definition}

\begin{result}[Integrated Brownian motion is Gaussian]
\label{res:ibm}
$Z$ is a Gaussian process with mean $0$ and, for $s\le t$,
\[
  \Cov\big(Z(s),Z(t)\big)=s^2\Big(\frac t2-\frac s6\Big);\qquad\text{in particular }\Var Z(t)=\frac{t^3}{3}.
\]
\end{result}
\begin{proof}
An integral is a limit of sums, and sums of values of $W$ are normal.
\begin{steps}
  \item $Z(t)$ is a limit of Riemann sums $\sum_kW(s_k)\Delta s$, which are MVN; limits of normal vectors are normal, so $Z$ is Gaussian.
  \item Mean: $\E_PZ(t)=\int_0^t\E_PW(s)\,ds=0$.
  \item Covariance: swap $\E_P$ and the integrals,
  \[
    \begin{aligned}
    \E_P\big(Z(s)Z(t)\big)&=\int_0^s\!\!\int_0^t\min(y,u)\,du\,dy=\int_0^s\Big[\int_0^yu\,du+\int_y^ty\,du\Big]dy\\
    &=\int_0^s\Big[\frac{y^2}{2}+y(t-y)\Big]dy=\int_0^s\Big[yt-\frac{y^2}{2}\Big]dy=\frac{ts^2}{2}-\frac{s^3}{6}. \qedhere
    \end{aligned}
  \]
\end{steps}
\end{proof}

\begin{pic}
\begin{tikzpicture}
\begin{axis}[width=10cm,height=4.6cm,xmin=0,xmax=1,xlabel={$t$},legend style={font=\scriptsize,at={(0.02,0.98)},anchor=north west},legend columns=2]
  \addplot[defc!50,thick] table[x=t,y=w1] {figs/ibm.dat}; \addlegendentry{$W(t)$}
  \addplot[defc,very thick] table[x=t,y=z1] {figs/ibm.dat}; \addlegendentry{$Z(t)=\int_0^tW$}
\end{axis}
\end{tikzpicture}
\picnote{Integrating smooths the path: $Z$ has a slope at every time (its slope is $W$), unlike $W$ itself.}
\end{pic}

%=====================================================================
\formulasheet{fs:gaussian}

\begin{mastercard}[label=card:mvn]{Jointly normal variables and Gaussian processes}
For jointly normal $X_1,\dots,X_n$ with means $\mu_i$ and covariances $C_{ij}=\Cov(X_i,X_j)$, and numbers $t_1,\dots,t_n$:
\[
  \boxed{\begin{gathered}\sum_it_iX_i\sim N\Big(\sum_it_i\mu_i,\ \sum_i\sum_jt_it_jC_{ij}\Big),\qquad \E_P\,e^{\sum_it_iX_i}=e^{\sum_it_i\mu_i+\frac12\sum_{i,j}t_it_jC_{ij}};\\
  X\given(Y=y)\sim N\Big(\mu_X+\frac{C_{XY}}{C_{YY}}(y-\mu_Y),\ C_{XX}-\frac{C_{XY}^2}{C_{YY}}\Big);\\
  \text{for }x_0+\mu t+\sigma W_t:\quad C(s,t)=\sigma^2\min(s,t).\end{gathered}}
\]
Values of $W$ at fixed times, combined with non-random coefficients, are always jointly normal. Uncorrelated jointly normal variables are independent.
\begin{plugrows}
  \plugrow{A combination of values of $W$}{$C_{ij}=\min(t_i,t_j)$}{normal, mean $0$, variance $\sum t_i^2\,t_i+2\sum_{i<j}t_it_j\min$}{count each cross term twice}{$4W_1+6W_3\sim N(0,172)$; $W_1-W_2+3W_4\sim N(0,31)$; $B(s)+B(t)\sim N(0,3s+t)$}
  \plugrow{Exponentials of $W$}{second formula}{$\E_Pe^{\sum t_iW_{t_i}}$; $\Cov(e^{X},e^{Y})=\E_Pe^{X+Y}-\E_Pe^X\E_Pe^Y$}{$X,Y$ jointly normal}{$\E_Pe^{W_1+W_2}=e^{5/2}$; $\Cov(e^{W_1},e^{W_2})=e^{3/2}(e-1)\approx7.70$}
  \plugrow{Products and odd moments}{condition on earlier times, or use symmetry}{odd-order products of mean-$0$ jointly normal variables have mean $0$}{means $0$}{$\E_P[B(t_1)B(t_2)B(t_3)]=0$}
  \plugrow{Conditioning in a general normal pair}{third formula}{normal law of $X$ given $Y$}{for $x_0+\mu t+\sigma W$ use Card~\ref{card:Ylaw} directly}{$(X,Y)$ with $C_{XX}=4$, $C_{YY}=2$, $C_{XY}=1$, means $0$: $X\given Y=2\sim N(1,3.5)$}
  \plugrow{Is it a Gaussian process? Is it a Wiener process?}{linear in values of $W$ with non-random coefficients $\Rightarrow$ Gaussian; then compare mean and $C(s,t)$ with $0$ and $\min(s,t)$, and check $X(0)=0$, continuity}{Wiener iff all four hold}{random coefficients or $W^2$, $e^W$ are not Gaussian}{$\frac1cW(c^2t)$ and $tW(1/t)$: Wiener; $e^{-t}W(e^{2t})$: Gaussian, $C=e^{-|t-s|}$, not Wiener}
  \plugrow{Brownian bridge}{$Z(t)=W(t)-tW(1)$ (or $W$ given $W(1)=0$)}{mean $0$, $C(s,t)=s(1-t)$ for $s\le t$}{$0\le s\le t\le1$}{$\Var Z(\frac12)=\frac14$}
  \plugrow{Integrated Brownian motion}{$Z(t)=\int_0^tW(u)\,du$; swap $\E_P$ and the integrals}{mean $0$, $C(s,t)=s^2\big(\frac t2-\frac s6\big)$ for $s\le t$; $\Var Z(t)=\frac{t^3}{3}$}{$s\le t$}{$\Var\int_0^1W=\frac13$}
  \plugrow{A joint density}{density of the first $\times$ conditional density of the second}{e.g.\ $(2W_3,W_5)$: $2W_3\sim N(0,12)$, $W_5\given2W_3=x\sim N(\frac x2,2)$}{the second given the first must be normal (true for BM)}{$f(x,y)=\frac{1}{4\pi\sqrt6}\exp\{-\frac{5x^2}{48}+\frac{xy}4-\frac{y^2}4\}$}
\end{plugrows}
\end{mastercard}
